\documentclass{article}
\usepackage[utf8]{inputenc}
\usepackage{amsmath}
\usepackage{amsfonts}
\usepackage{enumitem}

\begin{document}

\section*{Full Questions}
\begin{enumerate}
    \item Diketahui $a$ dan $b$ dua bilangan real nonnegatif yang memenuhi $a+2b=15$. Nilai terbesar yang mungkin dari $3a+5b$ adalah . . . .
    \begin{enumerate}[label={(\Alph*)}]
        \item 49 \item 45 \item 51 \item 41
    \end{enumerate}

    \item Dua benteng berbeda diletakkan pada sembarang petak di papan catur $8\times 8$. Dua buah benteng dikatakan \textit{saling menyerang} jika berada pada baris atau kolom yang sama. Peluang kedua benteng tidak saling menyerang adalah . . . .
    \begin{enumerate}[label={(\Alph*)}]
        \item $\dfrac{49}{64}$ \item $\dfrac{7}{9}$ \item $\dfrac{3}{4}$ \item $\dfrac{16}{21}$
    \end{enumerate}

    \item Diberikan persegi panjang $ABCD$ dengan luas $300$ satuan luas. Titik $E$ dan $F$ berturut-turut terletak pada sisi $AB$ dan $BC$ sedemikian sehingga berlaku perbandingan panjang $AE:EB=1:2$ dan $BF:FC=1:1$. Luas dari segiempat $DEFC$ adalah . . . satuan luas.
    \begin{enumerate}[label={(\Alph*)}]
        \item 150 \item 200 \item 160 \item 240
    \end{enumerate}

    \item Suatu bilangan asli disebut $bagus$ jika dapat dinyatakan sebagai $5a+6$, $6b+13$, dan $7c+1$ sekaligus untuk suatu bilangan bulat $a$, $b$, dan $c$. Perhatikan bahwa $1$ merupakan bilangan bagus karena
    $$1 = 5\times (-1)+6 = 6\times (-2) +13 = 7\times 0+1.$$
    Misalkan $N$ adalah bilangan asli terkecil selanjutnya yang bagus. Jumlah digit-digit dari $N$ adalah . . . .
    \begin{enumerate}[label={(\Alph*)}]
        \item 3 \item 4 \item 5 \item 6
    \end{enumerate}

    \item Dalam lima hari, koperasi $A$ berhasil menjual sebanyak $2,3,5,8$, dan $12$ seragam sekolah secara berturut-turut. Ketua koperasi ingin mengambil $2$ bilangan dari kelima bilangan tersebut, sehingga rata-rata dari dua bilangan yang dipilih lebih besar daripada rata-rata penjualan koperasi tersebut selama $5$ hari. Banyaknya cara yang mungkin dilakukan olehnya adalah . . . .
    \begin{enumerate}[label={(\Alph*)}]
        \item 2 \item 3 \item 4 \item 5
    \end{enumerate}

    \item Diketahui bahwa $x^2+px+q=0$ mempunyai $2$ akar bulat positif. Jika $p+q=4,$ maka nilai dari $q$ adalah . . . .
    \begin{enumerate}[label={(\Alph*)}]
        \item 8 \item 10 \item 12 \item 14
    \end{enumerate}

    \item Jin-ho mempunyai $64$ kubus berukuran $1\times 1\times 1$. Dari sekumpulan kubusnya, Jin-ho hendak menyusun seluruh kubus miliknya sehingga menjadi sebuah balok. Dua balok dikatakan berbeda jika salah satu balok tidak dapat dihasilkan dari refleksi, rotasi, ataupun translasi balok yang lainnya. Banyaknya macam balok berbeda yang bisa disusun oleh Jin-ho adalah . . . .
    \begin{enumerate}[label={(\Alph*)}]
        \item 4 \item 5 \item 6 \item 7
    \end{enumerate}

    \item Diberikan trapesium sama kaki $ABCD$ dengan $AB$ sejajar $CD$. Diketahui $\angle ACB=\angle ADB=30^\circ$ dan panjang $DA=AB=BC=1$. Luas dari $ABCD$ adalah . . . .
    \begin{enumerate}[label={(\Alph*)}]
        \item $\sqrt{3}$ \item $\dfrac{\sqrt{3}}{2}$ \item $\dfrac{3\sqrt{3}}{4}$ \item $\dfrac{\sqrt{3}}{4}$
    \end{enumerate}

    \item Banyaknya pasangan bilangan prima terurut $(p,q)$ sehingga $24-p-q$ juga bilangan prima adalah . . . .
    \begin{enumerate}[label={(\Alph*)}]
        \item 11 \item 13 \item 15 \item 17
    \end{enumerate}

    \item Andi memiliki $4$ koin yang memiliki $2$ sisi: sisi angka dan sisi gambar. Andi ingin mengetos keempat koin tersebut. Diketahui bahwa jika satu koin ditos, peluang munculnya sisi angka adalah $\frac{2}{5}$. Peluang bahwa agar terdapat minimal dua koin yang menunjukkan sisi gambar dalam pengetosan keempat koin tersebut adalah . . . .
    \begin{enumerate}[label={(\Alph*)}]
        \item $\dfrac{513}{625}$ \item $\dfrac{112}{625}$ \item $\dfrac{171}{625}$ \item $\dfrac{441}{625}$
    \end{enumerate}

    \item Diketahui $y+z$, $x+y$, $x+z$, dan $x+y+z$ merupakan empat suku pertama dari suatu barisan aritmetika di mana $x$, $y$, dan $z$ bilangan real. Jika suku kelima dari barisan tersebut adalah $168$, suku keenam dari barisan tersebut adalah . . . .
    \begin{enumerate}[label={(\Alph*)}]
        \item 180 \item 192 \item 204 \item 216
    \end{enumerate}

    \item Diberikan 100 koin yang semuanya menunjukkan sisi gambar. Setiap langkah dihitung dengan membalik sembarang 7 koin berbeda. Banyaknya langkah minimum sehingga semua koin menunjukkan sisi angka adalah . . . .
    \begin{enumerate}[label={(\Alph*)}]
        \item 15 \item 16 \item 17 \item 18
    \end{enumerate}

    \item Diberikan segidelapan beraturan $ABCDEFGH$ dengan panjang sisi $1$. Misalkan dua segmen $AC$ dan $BE$ saling berpotongan di titik $I$ serta dua segmen $EG$ dan $FA$ saling berpotongan di titik $J$. Jika luas dari $ABCDEFGH$ adalah $X$ dan luas $AIEJ$ adalah $Y$, nilai dari $\frac{X}{Y}$ adalah . . . .
    \begin{enumerate}[label={(\Alph*)}]
        \item $\sqrt{2}$ \item $\sqrt{2}-1$ \item $2\sqrt{2}$ \item $\sqrt{2}+1$
    \end{enumerate}

    \item Diketahui $a,b,c$ bilangan bulat tak negatif di mana $a\le 4$ dan $c>b>3$ yang memenuhi
    $$a+9b^2  + bc+6b= 125.$$
    Nilai terkecil dari $a+b+c$ adalah . . . .
    \begin{enumerate}[label={(\Alph*)}]
        \item 9 \item 13 \item 17 \item 22
    \end{enumerate}

    \item Ruby, Kana, dan Akane mengikuti lomba lari dengan $10$ peserta (termasuk mereka bertiga). Diketahui tidak ada dua pelari yang menyelesaikan pertandingan pada saat bersamaan. Peluang peringkat Ruby lebih rendah dari Kana dan Akane adalah . . . .
    \begin{enumerate}[label={(\Alph*)}]
        \item $\dfrac{1}{3}$ \item $\dfrac{1}{4}$ \item $\dfrac{1}{5}$ \item $\dfrac{1}{6}$
    \end{enumerate}

    \item Jika $x$ bilangan real, nilai terkecil dari
    $$\left (x^2-x+1\right )^2 -\left (x^2-x+1\right )+1$$
    yang mungkin adalah . . . .
    \begin{enumerate}[label={(\Alph*)}]
        \item $\dfrac{4}{5}$ \item $\dfrac{13}{16}$ \item $\dfrac{4}{7}$ \item $\dfrac{15}{22}$
    \end{enumerate}

    \item Seekor semut awalnya berada pada posisi $(0, 0)$ pada bidang kartesius. Setiap detiknya si semut berpindah tepat $1$ unit ke kanan, atau ke kiri, atau ke atas, atau ke bawah. Peluang bahwa setelah $4$ detik si semut berada di $(0, 0)$ adalah . . . .
    \begin{enumerate}[label={(\Alph*)}]
        \item $\dfrac{39}{256}$ \item $\dfrac{19}{128}$ \item $\dfrac{9}{64}$ \item $\dfrac{1}{4}$
    \end{enumerate}

    \item Pada segi-dua belas beraturan $ABCDEFGHIJKL$ berlaku panjang $AG=20$. Jarak titik $C$ ke garis $GL$ adalah . . . .
    \begin{enumerate}[label={(\Alph*)}]
        \item $10\sqrt{2}$ \item 10 \item $5\sqrt{6}$ \item 5
    \end{enumerate}

    \item Banyaknya bilangan asli kurang dari atau sama dengan $100$ yang jumlah semua faktor positifnya merupakan bilangan ganjil adalah . . . .
    \begin{enumerate}[label={(\Alph*)}]
        \item 16 \item 17 \item 18 \item 19
    \end{enumerate}

    \item Diketahui bahwa di dalam suatu kelas yang berisi $30$ murid, seorang guru matematika membagi kelas ini ke dalam 3 kelompok. Kelompok pertama berisi $a\ge 1$ murid yang semua nilai ulangannya $80$. Kelompok kedua berisi $b\ge 1$ murid yang semua nilai ulangannya $90$. Sedangkan kelompok ketiga berisi $c\ge 1$ murid yang semua nilai ulangannya $100$. Diketahui bahwa rata-rata nilai ulangan secara keseluruhan di kelas tersebut adalah $85$. Banyaknya tripel $(a,b,c)$ yang mungkin adalah . . . .
    \begin{enumerate}[label={(\Alph*)}]
        \item 6 \item 7 \item 8 \item 9
    \end{enumerate}

    % 1/113
    \item Diketahui
    $$ A =\frac{1}{1\times 2} + \frac{1}{3\times 4}+\frac{1}{4\times 5}+\cdots + \frac{1}{149\times 150} $$
    dan 
    $$ B =\frac{1}{76\times 150} + \frac{1}{77\times 149} + \frac{1}{78\times 148}+\cdots + \frac{1}{150\times 76}. $$
    Nilai dari $\frac{B}{A}$ adalah . . . .
    \begin{enumerate}[label={(\Alph*)}]
        \item $\dfrac{1}{111}$ \item $\dfrac{1}{113}$ \item $\dfrac{1}{115}$ \item $\dfrac{1}{117}$
    \end{enumerate}

    % 400
    \item Diberikan sebuah papan $8\times 8$ yang diisi sebuah bilangan di setiap kotaknya, dengan aturan bahwa kotak di baris ke-$i$ dan kolom ke-$j$ diisi $i\times j$. Kemudian, semua kotak yang memuat bilangan genap akan diwarnai hitam. Jumlah dari semua bilangan pada kotak diwarnai hitam adalah . . . .
    \begin{enumerate}[label={(\Alph*)}]
        \item 400 \item 720 \item 800 \item 1040
    \end{enumerate}

    % 78
    \item Di dalam persegi $ABCD$ terdapat dua titik $E$ dan $F$ yang memenuhi $\angle DEF=\angle EFB=60^\circ$. Jika panjang $DE=6$, $EF=10$, dan $FB=8$, maka luas dari persegi $ABCD$ adalah . . . .
    \begin{enumerate}[label={(\Alph*)}]
        \item 78 \item 84 \item 88 \item 62
    \end{enumerate}

    % 0
    \item Banyaknya pasangan bilangan asli $(a,b)$ yang memenuhi
    $$\frac{a}{b} + \frac{1-2b}{a} =4$$
    adalah . . . .
    \begin{enumerate}[label={(\Alph*)}]
        \item 0 \item 1 \item 2 \item 3
    \end{enumerate}

    % 253
    \item Banyaknya pasangan empat bilangan asli $(a,b,c,d)$ dengan $a<b<c<d$ sehingga himpunan $\{a,b,c,d\}$ memiliki median dan mean $24$ adalah . . . .
    \begin{enumerate}[label={(\Alph*)}]
        \item 253 \item 276 \item 506 \item 552
    \end{enumerate}

\end{enumerate}
\section*{Aljabar}
\begin{enumerate}[label={Soal \arabic*}]
    \item Diketahui $a$ dan $b$ dua bilangan real nonnegatif yang memenuhi $a+2b=15$. Nilai terbesar yang mungkin dari $3a+5b$ adalah . . . .
    \begin{enumerate}[label={(\Alph*)}]
        \item 49 \item 45 \item 51 \item 41
    \end{enumerate}


    \item Diketahui $y+z$, $x+y$, $x+z$, dan $x+y+z$ merupakan empat suku pertama dari suatu barisan aritmetika di mana $x$, $y$, dan $z$ bilangan real. Jika suku kelima dari barisan tersebut adalah $168$, suku keenam dari barisan tersebut adalah . . . .
    \begin{enumerate}[label={(\Alph*)}]
        \item 180 \item 192 \item 204 \item 216
    \end{enumerate}

    \item Diketahui $a,b,c$ bilangan bulat tak negatif di mana $a\le 4$ dan $c>b>3$ yang memenuhi
    $$a+9b^2  + bc+6b= 125.$$
    Nilai terkecil dari $a+b+c$ adalah . . . .
    \begin{enumerate}[label={(\Alph*)}]
        \item 9 \item 13 \item 17 \item 22
    \end{enumerate}

    \item Jika $x$ bilangan real, nilai terkecil dari
    $$\left (x^2-x+1\right )^2 -\left (x^2-x+1\right )+1$$
    yang mungkin adalah . . . .
    \begin{enumerate}[label={(\Alph*)}]
        \item $\dfrac{4}{5}$ \item $\dfrac{13}{16}$ \item $\dfrac{4}{7}$ \item $\dfrac{15}{22}$
    \end{enumerate}

    \item Diketahui
    $$ A =\frac{1}{1\times 2} + \frac{1}{3\times 4}+\frac{1}{4\times 5}+\cdots + \frac{1}{149\times 150} $$
    dan 
    $$ B =\frac{1}{76\times 150} + \frac{1}{77\times 149} + \frac{1}{78\times 148}+\cdots + \frac{1}{150\times 76}. $$
    Nilai dari $\frac{B}{A}$ adalah . . . .
    \begin{enumerate}[label={(\Alph*)}]
        \item $\dfrac{1}{111}$ \item $\dfrac{1}{113}$ \item $\dfrac{1}{115}$ \item $\dfrac{1}{117}$
    \end{enumerate}
\end{enumerate}

\section*{Teori Bilangan}
\begin{enumerate}[label={Soal \arabic*}]
    \setcounter{enumi}{7}
    \item Diketahui bahwa $x^2+px+q=0$ mempunyai $2$ akar bulat positif. Jika $p+q=4,$ maka nilai dari $q$ adalah . . . .
    \begin{enumerate}[label={(\Alph*)}]
        \item 8 \item 10 \item 12 \item 14
    \end{enumerate}

    \item Banyaknya pasangan bilangan asli $(a,b)$ yang memenuhi
    $$\frac{a}{b} + \frac{1-2b}{a} =4$$
    adalah . . . .
    \begin{enumerate}[label={(\Alph*)}]
        \item 0 \item 1 \item 2 \item 3
    \end{enumerate}
    
    \item Suatu bilangan asli disebut $bagus$ jika dapat dinyatakan sebagai $5a+6$, $6b+13$, dan $7c+1$ sekaligus untuk suatu bilangan bulat $a$, $b$, dan $c$. Perhatikan bahwa $1$ merupakan bilangan bagus karena
    $$1 = 5\times (-1)+6 = 6\times (-2) +13 = 7\times 0+1.$$
    Misalkan $N$ adalah bilangan asli terkecil selanjutnya yang bagus. Jumlah digit-digit dari $N$ adalah . . . .
    \begin{enumerate}[label={(\Alph*)}]
        \item 3 \item 4 \item 5 \item 6
    \end{enumerate}

    \item Banyaknya pasangan bilangan prima terurut $(p,q)$ sehingga $24-p-q$ juga bilangan prima adalah . . . .
    \begin{enumerate}[label={(\Alph*)}]
        \item 11 \item 13 \item 15 \item 17
    \end{enumerate}

    \item Banyaknya bilangan asli kurang dari atau sama dengan $100$ yang jumlah semua faktor positifnya merupakan bilangan ganjil adalah . . . .
    \begin{enumerate}[label={(\Alph*)}]
        \item 16 \item 17 \item 18 \item 19
    \end{enumerate}
\end{enumerate}

\section*{Geometri}
\begin{enumerate}[label={Soal \arabic*}]
    \setcounter{enumi}{10}
    \item Diberikan persegi panjang $ABCD$ dengan luas $300$ satuan luas. Titik $E$ dan $F$ berturut-turut terletak pada sisi $AB$ dan $BC$ sedemikian sehingga berlaku perbandingan panjang $AE:EB=1:2$ dan $BF:FC=1:1$. Luas dari segiempat $DEFC$ adalah . . . satuan luas.
    \begin{enumerate}[label={(\Alph*)}]
        \item 150 \item 200 \item 160 \item 240
    \end{enumerate}

    \item Diberikan trapesium sama kaki $ABCD$ dengan $AB$ sejajar $CD$. Diketahui $\angle ACB=\angle ADB=30^\circ$ dan panjang $DA=AB=BC=1$. Luas dari $ABCD$ adalah . . . .
    \begin{enumerate}[label={(\Alph*)}]
        \item $\sqrt{3}$ \item $\dfrac{\sqrt{3}}{2}$ \item $\dfrac{3\sqrt{3}}{4}$ \item $\dfrac{\sqrt{3}}{4}$
    \end{enumerate}

    \item Diberikan segidelapan beraturan $ABCDEFGH$ dengan panjang sisi $1$. Misalkan dua segmen $AC$ dan $BE$ saling berpotongan di titik $I$ serta dua segmen $EG$ dan $FA$ saling berpotongan di titik $J$. Jika luas dari $ABCDEFGH$ adalah $X$ dan luas $AIEJ$ adalah $Y$, nilai dari $\frac{X}{Y}$ adalah . . . .
    \begin{enumerate}[label={(\Alph*)}]
        \item $\sqrt{2}$ \item $\sqrt{2}-1$ \item $2\sqrt{2}$ \item $\sqrt{2}+1$
    \end{enumerate}

    \item Pada segi-dua belas beraturan $ABCDEFGHIJKL$ berlaku panjang $AG=20$. Jarak titik $C$ ke garis $GL$ adalah . . . .
    \begin{enumerate}[label={(\Alph*)}]
        \item $10\sqrt{2}$ \item 10 \item $5\sqrt{6}$ \item 5
    \end{enumerate}

    \item Di dalam persegi $ABCD$ terdapat dua titik $E$ dan $F$ yang memenuhi $\angle DEF=\angle EFB=60^\circ$. Jika panjang $DE=6$, $EF=10$, dan $FB=8$, maka luas dari persegi $ABCD$ adalah . . . .
    \begin{enumerate}[label={(\Alph*)}]
        \item 78 \item 84 \item 88 \item 62
    \end{enumerate}
\end{enumerate}

\section*{Kombinatorika}
\begin{enumerate}[label={Soal \arabic*}]
    \setcounter{enumi}{15}
    \item Dua benteng berbeda diletakkan pada sembarang petak di papan catur $8\times 8$. Dua buah benteng dikatakan \textit{saling menyerang} jika berada pada baris atau kolom yang sama. Peluang kedua benteng tidak saling menyerang adalah . . . .
    \begin{enumerate}[label={(\Alph*)}]
        \item $\dfrac{49}{64}$ \item $\dfrac{7}{9}$ \item $\dfrac{3}{4}$ \item $\dfrac{16}{21}$
    \end{enumerate}

    \item Dalam lima hari, koperasi $A$ berhasil menjual sebanyak $2,3,5,8$, dan $12$ seragam sekolah secara berturut-turut. Ketua koperasi ingin mengambil $2$ bilangan dari kelima bilangan tersebut, sehingga rata-rata dari dua bilangan yang dipilih lebih besar daripada rata-rata penjualan koperasi tersebut selama $5$ hari. Banyaknya cara yang mungkin dilakukan olehnya adalah . . . .
    \begin{enumerate}[label={(\Alph*)}]
        \item 2 \item 3 \item 4 \item 5
    \end{enumerate}

    \item Jin-ho mempunyai $64$ kubus berukuran $1\times 1\times 1$. Dari sekumpulan kubusnya, Jin-ho hendak menyusun seluruh kubus miliknya sehingga menjadi sebuah balok. Dua balok dikatakan berbeda jika salah satu balok tidak dapat dihasilkan dari refleksi, rotasi, ataupun translasi balok yang lainnya. Banyaknya macam balok berbeda yang bisa disusun oleh Jin-ho adalah . . . .
    \begin{enumerate}[label={(\Alph*)}]
        \item 4 \item 5 \item 6 \item 7
    \end{enumerate}

    \item Andi memiliki $4$ koin yang memiliki $2$ sisi: sisi angka dan sisi gambar. Andi ingin mengetos keempat koin tersebut. Diketahui bahwa jika satu koin ditos, peluang munculnya sisi angka adalah $\frac{2}{5}$. Peluang bahwa agar terdapat minimal dua koin yang menunjukkan sisi gambar dalam pengetosan keempat koin tersebut adalah . . . .
    \begin{enumerate}[label={(\Alph*)}]
        \item $\dfrac{513}{625}$ \item $\dfrac{112}{625}$ \item $\dfrac{171}{625}$ \item $\dfrac{441}{625}$
    \end{enumerate}

    \item Diberikan 100 koin yang semuanya menunjukkan sisi gambar. Setiap langkah dihitung dengan membalik sembarang 7 koin berbeda. Banyaknya langkah minimum sehingga semua koin menunjukkan sisi angka adalah . . . .
    \begin{enumerate}[label={(\Alph*)}]
        \item 15 \item 16 \item 17 \item 18
    \end{enumerate}

    \item Ruby, Kana, dan Akane mengikuti lomba lari dengan $10$ peserta (termasuk mereka bertiga). Diketahui tidak ada dua pelari yang menyelesaikan pertandingan pada saat bersamaan. Peluang peringkat Ruby lebih rendah dari Kana dan Akane adalah . . . .
    \begin{enumerate}[label={(\Alph*)}]
        \item $\dfrac{1}{3}$ \item $\dfrac{1}{4}$ \item $\dfrac{1}{5}$ \item $\dfrac{1}{6}$
    \end{enumerate}

    \item Seekor semut awalnya berada pada posisi $(0, 0)$ pada bidang kartesius. Setiap detiknya si semut berpindah tepat $1$ unit ke kanan, atau ke kiri, atau ke atas, atau ke bawah. Peluang bahwa setelah $4$ detik si semut berada di $(0, 0)$ adalah . . . .
    \begin{enumerate}[label={(\Alph*)}]
        \item $\dfrac{39}{256}$ \item $\dfrac{19}{128}$ \item $\dfrac{9}{64}$ \item $\dfrac{1}{4}$
    \end{enumerate}

    \item Diketahui bahwa di dalam suatu kelas yang berisi $30$ murid, seorang guru matematika membagi kelas ini ke dalam 3 kelompok. Kelompok pertama berisi $a\ge 1$ murid yang semua nilai ulangannya $80$. Kelompok kedua berisi $b\ge 1$ murid yang semua nilai ulangannya $90$. Sedangkan kelompok ketiga berisi $c\ge 1$ murid yang semua nilai ulangannya $100$. Diketahui bahwa rata-rata nilai ulangan secara keseluruhan di kelas tersebut adalah $85$. Banyaknya tripel $(a,b,c)$ yang mungkin adalah . . . .
    \begin{enumerate}[label={(\Alph*)}]
        \item 6 \item 7 \item 8 \item 9
    \end{enumerate}

    \item Diberikan sebuah papan $8\times 8$ yang diisi sebuah bilangan di setiap kotaknya, dengan aturan bahwa kotak di baris ke-$i$ dan kolom ke-$j$ diisi $i\times j$. Kemudian, semua kotak yang memuat bilangan genap akan diwarnai hitam. Jumlah dari semua bilangan pada kotak diwarnai hitam adalah . . . .
    \begin{enumerate}[label={(\Alph*)}]
        \item 400 \item 720 \item 800 \item 1040
    \end{enumerate}

    \item Banyaknya pasangan empat bilangan asli $(a,b,c,d)$ dengan $a<b<c<d$ sehingga himpunan $\{a,b,c,d\}$ memiliki median dan mean $24$ adalah . . . .
    \begin{enumerate}[label={(\Alph*)}]
        \item 253 \item 276 \item 506 \item 552
    \end{enumerate}
\end{enumerate}

\end{document}

\end{document}